% GENERATED FILE. Edit canonical lessons, then run npm run generate:books.
% canonical-source-hash: fnv1a64:70aeb5d138237ebb
% canonical-lessons: FR-C275-ecouler, FR-C275-ajourner, FR-C275-aboutir, FR-C275-contretemps, FR-C275-echeance, FR-R275-revision-268-275

\chapter{Le temps passe}
\label{ch:fr-le-temps-passe}
\hlchaptermodality{275}

\begin{chapteropening}
\textbf{By the end of this chapter:} \emph{I can say that time went by, that something was put off or came to a result, and name a hitch and a deadline.}

Complete the last lesson of chapter 275: I can say that time went by, that something was put off or came to a result, and name a hitch and a deadline.
\end{chapteropening}

\section[écouler]{écouler — to go by (as s'écouler, of time); to sell off (goods)}

\label{lesson:FR-C275-ecouler}



\begin{quote}
Before the new one: say the French for to precede, then the French for to last.
\end{quote}

\subsection*{You'll want to know: écouler}
\textbf{écouler} — \textquotedblleft{}to go by (as s'écouler, of time); to sell off (goods)\textquotedblright{}.

\begin{quote}
Trois ans se sont écoulés depuis son départ.
\end{quote}

\begin{quote}
\emph{Three years have gone by since he left.}
\end{quote}

Mostly used as \emph{s'écouler}: time goes by, and water flows out (\emph{l'eau s'écoule}). Without \emph{se}, it means to sell off goods: \emph{écouler un stock}.

\begin{cousinweb}
\emph{écouler} is built on \emph{couler} \textquotedblleft{}to flow\textquotedblright{}, which is inherited from Latin \emph{cōlāre} \textquotedblleft{}to strain, to filter\textquotedblright{} (Lewis \& Short).
\end{cousinweb}

\begin{grammarlens}[title={What you've built}]
One more word for telling, explaining and describing at B1.
\end{grammarlens}

\subsection*{Your turn}
\begin{itemize}
\raggedright
  \item \emph{Say it:} \emph{écouler}
  \item \emph{Say it:} \emph{écouler}, once more
  \item \emph{Say it:} a sentence of your own with \emph{écouler}
  \item \emph{Recall:} say the French for to precede, then the French for to last, then say \emph{écouler} again
\end{itemize}

\begin{tcolorbox}[breakable,colback=teal!4,colframe=teal!35!black,title={Before you move on}]
What does écouler mean? (\textquotedblleft{}To go by (as s'écouler, of time); to sell off (goods)\textquotedblright{}.) Say it once more.
\end{tcolorbox}
\section[ajourner]{ajourner — to postpone, to adjourn}

\label{lesson:FR-C275-ajourner}



\begin{quote}
Before the new one: say the French for to last, then the French for to go by.
\end{quote}

\subsection*{You'll want to know: ajourner}
\textbf{ajourner} — \textquotedblleft{}to postpone, to adjourn\textquotedblright{}.

\begin{quote}
Le match a été ajourné parce qu'il pleuvait.
\end{quote}

\begin{quote}
\emph{The match was postponed because it was raining.}
\end{quote}

Formal; in everyday speech, \emph{remettre à plus tard}.

\begin{cousinweb}
\emph{ajourner} comes from Old French \emph{ajorner} (EtymDB), built on \emph{jor} \textquotedblleft{}day\textquotedblright{}, from Latin \emph{diurnum} \textquotedblleft{}daily, a day's portion\textquotedblright{}: to put to another day.
\end{cousinweb}

\begin{grammarlens}[title={What you've built}]
One more word for telling, explaining and describing at B1.
\end{grammarlens}

\subsection*{Your turn}
\begin{itemize}
\raggedright
  \item \emph{Say it:} \emph{ajourner}
  \item \emph{Say it:} \emph{ajourner}, once more
  \item \emph{Say it:} a sentence of your own with \emph{ajourner}
  \item \emph{Recall:} say the French for to last, then the French for to go by, then say \emph{ajourner} again
\end{itemize}

\begin{tcolorbox}[breakable,colback=teal!4,colframe=teal!35!black,title={Before you move on}]
What does ajourner mean? (\textquotedblleft{}To postpone, to adjourn\textquotedblright{}.) Say it once more.
\end{tcolorbox}
\section[aboutir]{aboutir — to succeed, to come to a result; to lead to}

\label{lesson:FR-C275-aboutir}



\begin{quote}
Before the new one: say the French for to go by, then the French for to postpone.
\end{quote}

\subsection*{You'll want to know: aboutir}
\textbf{aboutir} — \textquotedblleft{}to succeed, to come to a result; to lead to\textquotedblright{}.

\begin{quote}
Le projet a enfin abouti.
\end{quote}

\begin{quote}
\emph{The project finally succeeded.}
\end{quote}

\emph{Aboutir à} — to lead to: \emph{la rue aboutit à la place}.

\begin{cousinweb}
\emph{aboutir} is built on \emph{à} and \emph{bout} \textquotedblleft{}end\textquotedblright{}: to reach the end. \emph{Bout} comes from \emph{bouter} \textquotedblleft{}to push\textquotedblright{}, which is usually traced to a Frankish verb; the origin is not settled (EtymDB lists no parent).
\end{cousinweb}

\begin{grammarlens}[title={What you've built}]
One more word for telling, explaining and describing at B1.
\end{grammarlens}

\subsection*{Your turn}
\begin{itemize}
\raggedright
  \item \emph{Say it:} \emph{aboutir}
  \item \emph{Say it:} \emph{aboutir}, once more
  \item \emph{Say it:} a sentence of your own with \emph{aboutir}
  \item \emph{Recall:} say the French for to go by, then the French for to postpone, then say \emph{aboutir} again
\end{itemize}

\begin{tcolorbox}[breakable,colback=teal!4,colframe=teal!35!black,title={Before you move on}]
What does aboutir mean? (\textquotedblleft{}To succeed, to come to a result; to lead to\textquotedblright{}.) Say it once more.
\end{tcolorbox}
\section[le contretemps]{le contretemps — the hitch, the mishap}

\label{lesson:FR-C275-contretemps}



\begin{quote}
Before the new one: say the French for to postpone, then the French for to succeed.
\end{quote}

\subsection*{You'll want to know: le contretemps}
\textbf{le contretemps} — \textquotedblleft{}the hitch, the mishap\textquotedblright{}.

\begin{quote}
Un petit contretemps nous a mis en retard.
\end{quote}

\begin{quote}
\emph{A small hitch made us late.}
\end{quote}

\emph{À contretemps} — at the wrong moment, out of step.

\begin{cousinweb}
\emph{contretemps} joins \emph{contre} \textquotedblleft{}against\textquotedblright{}, from Latin \emph{contrā}, and \emph{temps} \textquotedblleft{}time\textquotedblright{}, from Latin \emph{tempus} (Lewis \& Short): something that goes against the timing.
\end{cousinweb}

\begin{grammarlens}[title={What you've built}]
One more word for telling, explaining and describing at B1.
\end{grammarlens}

\subsection*{Your turn}
\begin{itemize}
\raggedright
  \item \emph{Say it:} \emph{le contretemps}
  \item \emph{Say it:} \emph{le contretemps}, once more
  \item \emph{Say it:} a sentence of your own with \emph{le contretemps}
  \item \emph{Recall:} say the French for to postpone, then the French for to succeed, then say \emph{le contretemps} again
\end{itemize}

\begin{tcolorbox}[breakable,colback=teal!4,colframe=teal!35!black,title={Before you move on}]
What does le contretemps mean? (\textquotedblleft{}The hitch, the mishap\textquotedblright{}.) Say it once more.
\end{tcolorbox}
\section[l'échéance]{l'échéance — the deadline, the due date}

\label{lesson:FR-C275-echeance}



\begin{quote}
Before the new one: say the French for to succeed, then the French for the hitch.
\end{quote}

\subsection*{You'll want to know: l'échéance}
\textbf{l'échéance} — \textquotedblleft{}the deadline, the due date\textquotedblright{}.

\begin{quote}
L'échéance approche et le travail n'est pas fini.
\end{quote}

\begin{quote}
\emph{The deadline is getting close and the work is not finished.}
\end{quote}

\emph{À brève échéance} — in the short term.

\begin{cousinweb}
\emph{échéance} is built on \emph{échoir} \textquotedblleft{}to fall due\textquotedblright{}, which is inherited from a Latin compound of \emph{ex} \textquotedblleft{}out\textquotedblright{} and \emph{cadere} \textquotedblleft{}to fall\textquotedblright{} (Lewis \& Short).
\end{cousinweb}

\begin{grammarlens}[title={What you've built}]
One more word for telling, explaining and describing at B1.
\end{grammarlens}

\subsection*{Your turn}
\begin{itemize}
\raggedright
  \item \emph{Say it:} \emph{l'échéance}
  \item \emph{Say it:} \emph{l'échéance}, once more
  \item \emph{Say it:} a sentence of your own with \emph{l'échéance}
  \item \emph{Recall:} say the French for to succeed, then the French for the hitch, then say \emph{l'échéance} again
\end{itemize}

\begin{tcolorbox}[breakable,colback=teal!4,colframe=teal!35!black,title={Before you move on}]
What does l'échéance mean? (\textquotedblleft{}The deadline, the due date\textquotedblright{}.) Say it once more.
\end{tcolorbox}
\section[(dialogue)]{Retour sur les faits}

\label{lesson:FR-R275-revision-268-275}



\begin{quote}
Say the French for first of all and in the first place.
\end{quote}

\subsection*{Your turn}
Say these aloud:

\begin{itemize}
\raggedright
  \item first of all, in the first place, afterwards, lastly, to finish
  \item when, as soon as, at the moment when, at once, while
  \item all of a sudden, suddenly, abruptly, at the time, straight after
  \item little by little, gradually, as one goes along, in the long run, sooner or later
  \item long ago, not long ago, before, until then, since then
  \item the day before, the next day, two days later, lately, soon
  \item to occur, to unfold, to take place, to precede, to last
  \item to go by, to postpone, to succeed, the hitch, the deadline
\end{itemize}

\begin{tcolorbox}[breakable,colback=teal!4,colframe=teal!35!black,title={Before you move on}]
First of all? (\textbf{tout d'abord}.) In the first place? (\textbf{en premier lieu}.) Afterwards? (\textbf{par la suite}.) Lastly? (\textbf{en dernier lieu}.) To finish? (\textbf{pour finir}.) When? (\textbf{lorsque}.) As soon as? (\textbf{dès que}.) At the moment when? (\textbf{au moment où}.) At once? (\textbf{aussitôt}.) While? (\textbf{tandis que}.) All of a sudden? (\textbf{tout à coup}.) Suddenly? (\textbf{subitement}.) Abruptly? (\textbf{brusquement}.) At the time? (\textbf{sur le coup}.) Straight after? (\textbf{dans la foulée}.) Little by little? (\textbf{peu à peu}.) Gradually? (\textbf{progressivement}.) As one goes along? (\textbf{au fur et à mesure}.) In the long run? (\textbf{à la longue}.) Sooner or later? (\textbf{tôt ou tard}.) Long ago? (\textbf{jadis}.) Not long ago? (\textbf{naguère}.) Before? (\textbf{auparavant}.) Until then? (\textbf{jusqu'alors}.) Since then? (\textbf{depuis lors}.) The day before? (\textbf{la veille}.) The next day? (\textbf{le lendemain}.) Two days later? (\textbf{le surlendemain}.) Lately? (\textbf{dernièrement}.) Soon? (\textbf{prochainement}.) To occur? (\textbf{survenir}.) To unfold? (\textbf{dérouler}.) To take place? (\textbf{avoir lieu}.) To precede? (\textbf{précéder}.) To last? (\textbf{durer}.) To go by? (\textbf{écouler}.) To postpone? (\textbf{ajourner}.) To succeed? (\textbf{aboutir}.) The hitch? (\textbf{le contretemps}.) The deadline? (\textbf{l'échéance}.)
\end{tcolorbox}
